% Structured abstract — JAMIA convention.
%
% Required subsections (per JAMIA General Instructions for Authors):
%   Objective / Materials and Methods / Results / Discussion / Conclusion
% Word limit: 250 prose words (subsection labels do not count).
%
% Subsection labels are rendered as BOLD CAPS per JAMIA's heading
% hierarchy spec ("BOLD CAPS, bold lower case, Plain text, Italics").

\begin{center}
{\Large\bfseries ABSTRACT\par}
\end{center}

\noindent\textbf{OBJECTIVE.}\\
To evaluate the most recent frontier large language models (LLMs)
on medical concept normalization across both safety (fabrication
of non-existent codes vs.\ abstention) and deployment cost (cost-
per-correct), and to identify the (model, prompting-mode)
configurations that are Pareto-optimal for production deployment.

\medskip
\noindent\textbf{MATERIALS AND METHODS.}\\
We evaluate eight frontier LLMs under three prompting modes
(zero-shot, constrained, retrieval-augmented) against
string-matching and frozen sentence-transformer baselines.
Inputs are Synthea-generated FHIR Conditions under four
degradation modes. A six-way outcome taxonomy separates
fabrication from abstention.

\medskip
\noindent\textbf{RESULTS.}\\
On 125 unique Synthea Conditions across 4 degradation modes,
fabrication ranged from 0 to 5.6\% across flagship configurations
and from 0 to 12.8\% across sub-flagship configurations.
Constrained prompting eliminated fabrication for every Anthropic
and OpenAI model. Abstention proved more consequential for some
models: Gemini 2.5 Pro returned an empty prediction on 88\% of
zero-shot D4 inputs. Cost-per-correct varied by a factor of 270
across the matrix; GPT-4o-mini in constrained mode achieved the
lowest cost at \$0.0003 per exact match.

\medskip
\noindent\textbf{DISCUSSION.}\\
Hallucination appears more controllable than prior literature
suggested at this scale. Larger cohorts with more realistic
clinical text are needed to establish the true effect size.

\medskip
\noindent\textbf{CONCLUSION.}\\
Zero-shot LLM prompting is unsafe for medical concept
normalization; constrained or retrieval-augmented grounding is
required to eliminate fabrication. Among grounded configurations,
cost-per-correct drives deployment selection: GPT-4o-mini
constrained leads at 94.2\% top-1 and \$0.0003 per correct, with
Claude Haiku 4.5 constrained as the higher-accuracy alternative
(96.8\%, \$0.0044). Both dominate the largest frontier model
tested (Opus 4.7, \$0.0133 per correct).
